\textbf{I would not approve an equity investment on the case presented.
I would consider a small, milestone-funded diligence option, not a
venture-scale commitment.} The documents make a plausible public-funding
and laboratory-operating case. They do not establish a company-level
mechanism for capturing a durable share of the value created. Open
science can be socially valuable and a poor investment in the entity
that pays to produce it. Conversely, openness does not by itself make a
company uninvestable.

My recomputation finds \textbf{no reliable positive association between
openness and the reported value-to-funding proxy}. Among 118
identifiable commercial, non-model entities, only 37 have both an
admissible value and positive recorded funding. Scores 2-4 currently
have a median proxy of \textbf{7.54x} (n=23), versus \textbf{11.49x} for
scores 0-1 (n=14): median ratio \textbf{0.66}, 95\% bootstrap interval
\textbf{0.35-1.30}. At founding, the corresponding medians are
\textbf{7.90x} (n=25) and \textbf{9.69x} (n=12), ratio \textbf{0.82}
(0.50-1.53). These are neither evidence of an openness premium nor proof
that closed strategies win.

The dataset is a curated casebook, not an investable-company census. It
contains only \textbf{one explicit shutdown and one bankruptcy}, mixes
projects with companies, and sometimes puts revenue or grants in the
value column. After defensible filtering, there are \textbf{no currently
level-4 commercial entities}, and the category literally called physical
science has \textbf{zero usable values or funding multiples}. The
relevant lab-automation/AI-for-science group supplies just four
multiples, only one for a currently open-core company. It cannot
underwrite an alloy-discovery investment.

\textbf{The decision should turn on paid customer demand, reproducible
physical performance, fully burdened unit economics, an enforceable
IP/data boundary and a financing-to-exit plan.} Open drivers can be
useful distribution infrastructure. They are not a demonstrated moat.
